\documentclass{situationist-preprint}

\title{Why There Is Science}
\msid{SCI-18-b}
\msfamily{SCI-18}
\msdiscipline{Philosophy of Science}
\mstype{Companion manuscript}
\mskeywords{limits of explanation; enabling conditions; past hypothesis; anthropic reasoning; existence of science}

\begin{document}
\maketitle

\begin{abstract}
SCI-18-a argues that science cannot explain the existence of the conditions that make science possible. We apply the result to the question why there is science at all. The question divides into two: why the enabling conditions of science obtain, and why, given them, there are inquirers who practice science. Science can answer the second, and has largely done so. It cannot answer the first. We show where the division falls for the low-entropy initial state of the universe, which TSM-05-a counts among the enabling conditions because it makes records reliable, and we argue that anthropic reasoning and cosmological models of the origin of low entropy do not cross the line. The existence of science is partly explained by science, and the part that is not is exactly the part SCI-18-a identifies.
\end{abstract}

\section{The question}

Why is there science? The question can be heard in two ways. Heard as a question about human history, it asks why certain animals came to test hypotheses systematically, and it has scientific and historical answers. Heard as a question about the world, it asks why there is anything for science to be about, regularities stable enough to discover and records reliable enough to learn from. SCI-18-a argues, by its Theorem~3.1, that science cannot explain the existence of the enabling conditions of science, which by its Definition~2.2 include the existence of a concrete world, the holding of laws of nature, and the past hypothesis. The aim of this manuscript is to locate exactly where the limit falls when the question is why there is science.

\section{The division}

Let the \emph{enabling conditions} be as in Definition~2.2 of SCI-18-a, and let the \emph{practice conditions} be the further facts, about organisms, cognition and institutions, without which science would be possible but not actual.

\begin{proposition}\label{p:div}
That there is science is the conjunction of the enabling conditions and the practice conditions. The practice conditions admit scientific explanation relative to the enabling conditions; the enabling conditions do not admit scientific explanation.
\end{proposition}

\begin{proof}
The first claim is the definition of the two classes: science is actual just in case it is possible and the facts that make it practiced obtain. For the second, the practice conditions are facts about particular concrete systems: the evolution of organisms able to represent regularities, the development of writing and measurement, the formation of institutions of error-correction of the kind SCI-01 identifies with science. Each is a fact for which a scientific explanation in the sense of Definition~2.1 of SCI-18-a can cite laws and earlier conditions, and biology, psychology and history supply such explanations. For the third, Theorem~3.1 of SCI-18-a applies directly.
\end{proof}

The answer to the question why there is science is therefore partly scientific. It is scientific to the extent that it explains how inquirers arose in a world of the right kind, and it falls silent exactly where the question becomes why the world is of that kind.

\section{The past hypothesis}

The case of the past hypothesis shows that the line does not fall where one might expect. TSM-05-a argues that time has a direction because the universe began in a state of very low entropy, and it counts the past hypothesis among the conditions that make science possible, because it is what makes records, and therefore evidence about the past, reliable \citep{Albert2000}. Without it, the inference from a present record to the event it records would be no better than the inference from the record to a fluctuation that produced it. Every experiment relies on records of its own earlier stages; every confirmation relies on the reliability of reports. The past hypothesis is thus not one condition among the objects of science but a condition of the evidence by which science establishes anything.

\begin{corollary}\label{c:ph}
No scientific explanation of the past hypothesis is available that does not cite a further condition of the same kind, namely a low-entropy or otherwise special state of an earlier or larger system.
\end{corollary}

\begin{proof}
This is Corollary~3.2 of SCI-18-a. A scientific explanation of the past hypothesis would cite laws and a concrete condition; to explain why entropy was low, the condition must be one from which low entropy follows, and the reliability of the evidence for that condition would depend on the condition itself.
\end{proof}

Corollary~\ref{c:ph} agrees with TSM-05-a, which argues that the past hypothesis requires no further explanation: on the Humean account of laws it may be an axiom of the best system, and asking why it holds is asking why the mosaic is as it is. SCI-18-a and TSM-05-a reach compatible conclusions from different directions. The first says that science cannot supply an explanation; the second says that none is owed.

\section{Objections and replies}

\begin{objection}
Anthropic reasoning explains why the enabling conditions obtain: only in a world with stable laws and low initial entropy could there be observers to ask the question \citep{Carter1974}.
\end{objection}

\begin{reply}
Anthropic reasoning explains why observers find the enabling conditions satisfied, given that there are observers. It is a selection effect on what is observed, not an explanation of what exists. It accounts for a practice condition, our being located where the enabling conditions hold, and leaves the existence of the enabling conditions where it was. If supplemented by a multiverse in which they hold in some regions and not others, the explanation cites the laws and initial state of the multiverse, which are enabling conditions of the same kinds.
\end{reply}

\begin{objection}
Cosmological models in which our low-entropy beginning arises from a prior state, by inflation or by the nucleation of a universe from a larger spacetime, are scientific explanations of the past hypothesis.
\end{objection}

\begin{reply}
They explain our past hypothesis by a condition of the prior state, and that condition must itself be special if the explanation is to work: the measure of initial states that yield inflation of the required kind is itself a question about initial entropy \citep{Penrose1989}. Such models may be correct, and they would push the past hypothesis back. They do not eliminate it, which is what Corollary~\ref{c:ph} says.
\end{reply}

\begin{objection}
If science cannot explain why there is science, the existence of science is a brute fact, and inquiry rests on something arbitrary.
\end{objection}

\begin{reply}
SCI-18-a leaves open whether the enabling conditions have explanations of another kind. One of them is addressed elsewhere: MET-01 argues that there is something rather than nothing because the empty world is not a stable possibility, an explanation that is metaphysical rather than scientific, since it cites no laws and no concrete conditions. Whatever the status of the others, a condition that cannot be explained scientifically is not thereby arbitrary in any sense that would undermine the inquiry it makes possible. The reliability of records is not diminished by our inability to explain why entropy was once low.
\end{reply}

\section{Conclusion}

There is science because there are inquirers, and science explains how there came to be inquirers. There is science also because the world contains stable regularities and an initial state that makes evidence about the past reliable, and science cannot explain this, since every explanation it could offer presupposes conditions of the same kinds. The question why there is science has a scientific answer up to exactly the point SCI-18-a identifies.

\begin{thebibliography}{99}\small
\bibitem[Albert(2000)]{Albert2000} Albert, D.~Z. (2000). \emph{Time and Chance}. Harvard University Press.
\bibitem[Carter(1974)]{Carter1974} Carter, B. (1974). Large number coincidences and the anthropic principle in cosmology. In M.~S. Longair (ed.), \emph{Confrontation of Cosmological Theories with Observational Data}, IAU Symposium 63, 291--298. Reidel.
\bibitem[Penrose(1989)]{Penrose1989} Penrose, R. (1989). \emph{The Emperor's New Mind}. Oxford University Press.
\end{thebibliography}

\end{document}
